% ============================================================
% Verification Tables + Prose (Termination + Safety)
% F^{V,Delta,Sigma}_{Tendermint}  Delay-Attack Model
% ============================================================
%
% Required in preamble:
%   \usepackage{booktabs}
%   \usepackage{amsmath,amssymb}
%
% Insert in body:
%   % ============================================================
% Verification Tables + Prose (Termination + Safety)
% F^{V,Delta,Sigma}_{Tendermint}  Delay-Attack Model
% ============================================================
%
% Required in preamble:
%   \usepackage{booktabs}
%   \usepackage{amsmath,amssymb}
%
% Insert in body:
%   % ============================================================
% Verification Tables + Prose (Termination + Safety)
% F^{V,Delta,Sigma}_{Tendermint}  Delay-Attack Model
% ============================================================
%
% Required in preamble:
%   \usepackage{booktabs}
%   \usepackage{amsmath,amssymb}
%
% Insert in body:
%   % ============================================================
% Verification Tables + Prose (Termination + Safety)
% F^{V,Delta,Sigma}_{Tendermint}  Delay-Attack Model
% ============================================================
%
% Required in preamble:
%   \usepackage{booktabs}
%   \usepackage{amsmath,amssymb}
%
% Insert in body:
%   \input{verification_tables}
% ============================================================

% ==================================================================
% 1) TERMINATION
% ==================================================================
\subsection{Termination Verification}

For termination, we only use two key notions:
\(\mathrm{InitiallyHonest}\) and \(\mathrm{NotYetTerminated}\).

\begin{align*}
\mathrm{InitiallyHonest}
&\triangleq
V \setminus V_{\mathrm{corrupted}} \\
\mathrm{NotYetTerminated}
&\triangleq
\neg\Big(\forall p\in\mathrm{InitiallyHonest},\; \mathrm{decision}_p(0)\neq \bot\Big)
\end{align*}

Interpretation:
\(\mathrm{InitiallyHonest}\) is the set of nodes not initially corrupted;
\(\mathrm{NotYetTerminated}\) means that at least one initially honest node
has not decided at height \(0\).
In our runs, a counterexample to \texttt{NotYetTerminated}
(i.e., RC\(=12\)) is treated as a termination witness.

For Table~\ref{tab:termination}, each column has the following meaning:
\textbf{Run} is the experiment identifier,
$R_{\max}$ is the maximum round bound used in that run,
\textbf{length} is the simulation trace-length bound passed to Apalache,
\textbf{Time (s)} is wall-clock verification time in seconds,
and \textbf{Result} indicates whether a termination witness was found
(PASS corresponds to RC$=12$ in this setup).

\begin{table}[ht]
\centering
\caption{%
  Termination verification results for
  $\mathcal{F}^{V,\Delta,\Sigma}_{\mathrm{Tendermint}}$
  under bounded delay attack
  ($|V|=4$, $f=1$, $\Delta=\Sigma=1$, $H_{\max}=1$,
  \texttt{length}$=100$, \texttt{max-run}$=60$).
  RC$\,=\,12$ indicates that Apalache found a termination witness.
}
\label{tab:termination}
\small
\begin{tabular}{@{}ccrrc@{}}
\toprule
\textbf{Run} &
$R_{\max}$ &
\textbf{length} &
\textbf{Time\,(s)} &
\textbf{Result} \\
\midrule
TR-01 &  1 & 100 &    31.3 & PASS \\
TR-02 &  2 & 100 &    25.4 & PASS \\
TR-03 &  3 & 100 &    34.0 & PASS \\
TR-04 &  5 & 100 &   117.1 & PASS \\
TR-05 &  8 & 100 &    81.4 & PASS \\
TR-06 & 10 & 100 &  1097.1 & PASS \\
\bottomrule
\end{tabular}
\end{table}

% ==================================================================
% 2) SAFETY
% ==================================================================
\subsection{Safety Verification}

To verify safety of $\mathcal{F}^{V,\Delta,\Sigma}_{\mathrm{Tendermint}}$
under delay attacks, we use Apalache in bounded checking mode
(\texttt{apalache-mc check}) and evaluate the invariant
$\texttt{Agreement}$, requiring that no two initially honest validators
decide different values at the same height.

The safety property in standard logical form is:
\begin{align*}
\mathrm{Agreement}
&\triangleq
\forall h\in\{0,\dots,H_{\max}\},\; \forall p,q\in\mathrm{InitiallyHonest},
\\
&\qquad\Big(\mathrm{decision}_p(h)\neq\bot \wedge \mathrm{decision}_q(h)\neq\bot\Big)
\Rightarrow
\mathrm{decision}_p(h)=\mathrm{decision}_q(h)
\end{align*}
We fix the validator set $V=\{v_1,v_2,v_3,v_4\}$ with fault threshold $f=1$,
network delay bound $\Delta=1$, protocol attack tolerance $\Sigma=1$,
and block height bound $H_{\max}=1$.

The checker outcome is interpreted as follows:
RC$\,=\,0$ (NoError) means no Agreement counterexample exists within the
bounded horizon (here \texttt{length}$=5$), while RC$\,=\,12$ would indicate
that a violating trace was found.
Table~\ref{tab:safety} shows RC$=0$ for every tested $R_{\max}$,
so no safety violation is observed in this configuration.

Similar to other bounded symbolic checks, runtime is not strictly monotonic
in $R_{\max}$: solver effort depends on the structure of enabled/disabled
transitions and pruning opportunities, not only on the numeric round bound.

For Table~\ref{tab:safety}, the column semantics are:
\textbf{Run} denotes the experiment identifier,
$R_{\max}$ is the maximum round bound for that check,
\textbf{length} is the bounded model-checking horizon,
\textbf{Time (s)} is wall-clock runtime in seconds,
and \textbf{Result} reports the Agreement-check outcome
(PASS corresponds to RC$=0$, i.e., no counterexample found within bounds).

\begin{table}[ht]
\centering
\caption{%
  Safety verification results for
  $\mathcal{F}^{V,\Delta,\Sigma}_{\mathrm{Tendermint}}$
  under bounded delay attack
  ($|V|=4$, $f=1$, $\Delta=\Sigma=1$, $H_{\max}=1$,
  \texttt{length}$=5$).
  RC$\,=\,0$ indicates that no Agreement counterexample was found
  within the bounded computation.
}
\label{tab:safety}
\small
\begin{tabular}{@{}ccrrc@{}}
\toprule
\textbf{Run} &
$R_{\max}$ &
\textbf{length} &
\textbf{Time\,(s)} &
\textbf{Result} \\
\midrule
SR-01 &  1 & 5 & 61.6 & PASS \\
SR-02 &  2 & 5 & 60.6 & PASS \\
SR-03 &  3 & 5 & 57.7 & PASS \\
SR-04 &  5 & 5 & 48.1 & PASS \\
SR-05 &  8 & 5 & 58.3 & PASS \\
SR-06 & 10 & 5 & 44.9 & PASS \\
\bottomrule
\end{tabular}
\end{table}

% ============================================================

% ==================================================================
% 1) TERMINATION
% ==================================================================
\subsection{Termination Verification}

For termination, we only use two key notions:
\(\mathrm{InitiallyHonest}\) and \(\mathrm{NotYetTerminated}\).

\begin{align*}
\mathrm{InitiallyHonest}
&\triangleq
V \setminus V_{\mathrm{corrupted}} \\
\mathrm{NotYetTerminated}
&\triangleq
\neg\Big(\forall p\in\mathrm{InitiallyHonest},\; \mathrm{decision}_p(0)\neq \bot\Big)
\end{align*}

Interpretation:
\(\mathrm{InitiallyHonest}\) is the set of nodes not initially corrupted;
\(\mathrm{NotYetTerminated}\) means that at least one initially honest node
has not decided at height \(0\).
In our runs, a counterexample to \texttt{NotYetTerminated}
(i.e., RC\(=12\)) is treated as a termination witness.

For Table~\ref{tab:termination}, each column has the following meaning:
\textbf{Run} is the experiment identifier,
$R_{\max}$ is the maximum round bound used in that run,
\textbf{length} is the simulation trace-length bound passed to Apalache,
\textbf{Time (s)} is wall-clock verification time in seconds,
and \textbf{Result} indicates whether a termination witness was found
(PASS corresponds to RC$=12$ in this setup).

\begin{table}[ht]
\centering
\caption{%
  Termination verification results for
  $\mathcal{F}^{V,\Delta,\Sigma}_{\mathrm{Tendermint}}$
  under bounded delay attack
  ($|V|=4$, $f=1$, $\Delta=\Sigma=1$, $H_{\max}=1$,
  \texttt{length}$=100$, \texttt{max-run}$=60$).
  RC$\,=\,12$ indicates that Apalache found a termination witness.
}
\label{tab:termination}
\small
\begin{tabular}{@{}ccrrc@{}}
\toprule
\textbf{Run} &
$R_{\max}$ &
\textbf{length} &
\textbf{Time\,(s)} &
\textbf{Result} \\
\midrule
TR-01 &  1 & 100 &    31.3 & PASS \\
TR-02 &  2 & 100 &    25.4 & PASS \\
TR-03 &  3 & 100 &    34.0 & PASS \\
TR-04 &  5 & 100 &   117.1 & PASS \\
TR-05 &  8 & 100 &    81.4 & PASS \\
TR-06 & 10 & 100 &  1097.1 & PASS \\
\bottomrule
\end{tabular}
\end{table}

% ==================================================================
% 2) SAFETY
% ==================================================================
\subsection{Safety Verification}

To verify safety of $\mathcal{F}^{V,\Delta,\Sigma}_{\mathrm{Tendermint}}$
under delay attacks, we use Apalache in bounded checking mode
(\texttt{apalache-mc check}) and evaluate the invariant
$\texttt{Agreement}$, requiring that no two initially honest validators
decide different values at the same height.

The safety property in standard logical form is:
\begin{align*}
\mathrm{Agreement}
&\triangleq
\forall h\in\{0,\dots,H_{\max}\},\; \forall p,q\in\mathrm{InitiallyHonest},
\\
&\qquad\Big(\mathrm{decision}_p(h)\neq\bot \wedge \mathrm{decision}_q(h)\neq\bot\Big)
\Rightarrow
\mathrm{decision}_p(h)=\mathrm{decision}_q(h)
\end{align*}
We fix the validator set $V=\{v_1,v_2,v_3,v_4\}$ with fault threshold $f=1$,
network delay bound $\Delta=1$, protocol attack tolerance $\Sigma=1$,
and block height bound $H_{\max}=1$.

The checker outcome is interpreted as follows:
RC$\,=\,0$ (NoError) means no Agreement counterexample exists within the
bounded horizon (here \texttt{length}$=5$), while RC$\,=\,12$ would indicate
that a violating trace was found.
Table~\ref{tab:safety} shows RC$=0$ for every tested $R_{\max}$,
so no safety violation is observed in this configuration.

Similar to other bounded symbolic checks, runtime is not strictly monotonic
in $R_{\max}$: solver effort depends on the structure of enabled/disabled
transitions and pruning opportunities, not only on the numeric round bound.

For Table~\ref{tab:safety}, the column semantics are:
\textbf{Run} denotes the experiment identifier,
$R_{\max}$ is the maximum round bound for that check,
\textbf{length} is the bounded model-checking horizon,
\textbf{Time (s)} is wall-clock runtime in seconds,
and \textbf{Result} reports the Agreement-check outcome
(PASS corresponds to RC$=0$, i.e., no counterexample found within bounds).

\begin{table}[ht]
\centering
\caption{%
  Safety verification results for
  $\mathcal{F}^{V,\Delta,\Sigma}_{\mathrm{Tendermint}}$
  under bounded delay attack
  ($|V|=4$, $f=1$, $\Delta=\Sigma=1$, $H_{\max}=1$,
  \texttt{length}$=5$).
  RC$\,=\,0$ indicates that no Agreement counterexample was found
  within the bounded computation.
}
\label{tab:safety}
\small
\begin{tabular}{@{}ccrrc@{}}
\toprule
\textbf{Run} &
$R_{\max}$ &
\textbf{length} &
\textbf{Time\,(s)} &
\textbf{Result} \\
\midrule
SR-01 &  1 & 5 & 61.6 & PASS \\
SR-02 &  2 & 5 & 60.6 & PASS \\
SR-03 &  3 & 5 & 57.7 & PASS \\
SR-04 &  5 & 5 & 48.1 & PASS \\
SR-05 &  8 & 5 & 58.3 & PASS \\
SR-06 & 10 & 5 & 44.9 & PASS \\
\bottomrule
\end{tabular}
\end{table}

% ============================================================

% ==================================================================
% 1) TERMINATION
% ==================================================================
\subsection{Termination Verification}

For termination, we only use two key notions:
\(\mathrm{InitiallyHonest}\) and \(\mathrm{NotYetTerminated}\).

\begin{align*}
\mathrm{InitiallyHonest}
&\triangleq
V \setminus V_{\mathrm{corrupted}} \\
\mathrm{NotYetTerminated}
&\triangleq
\neg\Big(\forall p\in\mathrm{InitiallyHonest},\; \mathrm{decision}_p(0)\neq \bot\Big)
\end{align*}

Interpretation:
\(\mathrm{InitiallyHonest}\) is the set of nodes not initially corrupted;
\(\mathrm{NotYetTerminated}\) means that at least one initially honest node
has not decided at height \(0\).
In our runs, a counterexample to \texttt{NotYetTerminated}
(i.e., RC\(=12\)) is treated as a termination witness.

For Table~\ref{tab:termination}, each column has the following meaning:
\textbf{Run} is the experiment identifier,
$R_{\max}$ is the maximum round bound used in that run,
\textbf{length} is the simulation trace-length bound passed to Apalache,
\textbf{Time (s)} is wall-clock verification time in seconds,
and \textbf{Result} indicates whether a termination witness was found
(PASS corresponds to RC$=12$ in this setup).

\begin{table}[ht]
\centering
\caption{%
  Termination verification results for
  $\mathcal{F}^{V,\Delta,\Sigma}_{\mathrm{Tendermint}}$
  under bounded delay attack
  ($|V|=4$, $f=1$, $\Delta=\Sigma=1$, $H_{\max}=1$,
  \texttt{length}$=100$, \texttt{max-run}$=60$).
  RC$\,=\,12$ indicates that Apalache found a termination witness.
}
\label{tab:termination}
\small
\begin{tabular}{@{}ccrrc@{}}
\toprule
\textbf{Run} &
$R_{\max}$ &
\textbf{length} &
\textbf{Time\,(s)} &
\textbf{Result} \\
\midrule
TR-01 &  1 & 100 &    31.3 & PASS \\
TR-02 &  2 & 100 &    25.4 & PASS \\
TR-03 &  3 & 100 &    34.0 & PASS \\
TR-04 &  5 & 100 &   117.1 & PASS \\
TR-05 &  8 & 100 &    81.4 & PASS \\
TR-06 & 10 & 100 &  1097.1 & PASS \\
\bottomrule
\end{tabular}
\end{table}

% ==================================================================
% 2) SAFETY
% ==================================================================
\subsection{Safety Verification}

To verify safety of $\mathcal{F}^{V,\Delta,\Sigma}_{\mathrm{Tendermint}}$
under delay attacks, we use Apalache in bounded checking mode
(\texttt{apalache-mc check}) and evaluate the invariant
$\texttt{Agreement}$, requiring that no two initially honest validators
decide different values at the same height.

The safety property in standard logical form is:
\begin{align*}
\mathrm{Agreement}
&\triangleq
\forall h\in\{0,\dots,H_{\max}\},\; \forall p,q\in\mathrm{InitiallyHonest},
\\
&\qquad\Big(\mathrm{decision}_p(h)\neq\bot \wedge \mathrm{decision}_q(h)\neq\bot\Big)
\Rightarrow
\mathrm{decision}_p(h)=\mathrm{decision}_q(h)
\end{align*}
We fix the validator set $V=\{v_1,v_2,v_3,v_4\}$ with fault threshold $f=1$,
network delay bound $\Delta=1$, protocol attack tolerance $\Sigma=1$,
and block height bound $H_{\max}=1$.

The checker outcome is interpreted as follows:
RC$\,=\,0$ (NoError) means no Agreement counterexample exists within the
bounded horizon (here \texttt{length}$=5$), while RC$\,=\,12$ would indicate
that a violating trace was found.
Table~\ref{tab:safety} shows RC$=0$ for every tested $R_{\max}$,
so no safety violation is observed in this configuration.

Similar to other bounded symbolic checks, runtime is not strictly monotonic
in $R_{\max}$: solver effort depends on the structure of enabled/disabled
transitions and pruning opportunities, not only on the numeric round bound.

For Table~\ref{tab:safety}, the column semantics are:
\textbf{Run} denotes the experiment identifier,
$R_{\max}$ is the maximum round bound for that check,
\textbf{length} is the bounded model-checking horizon,
\textbf{Time (s)} is wall-clock runtime in seconds,
and \textbf{Result} reports the Agreement-check outcome
(PASS corresponds to RC$=0$, i.e., no counterexample found within bounds).

\begin{table}[ht]
\centering
\caption{%
  Safety verification results for
  $\mathcal{F}^{V,\Delta,\Sigma}_{\mathrm{Tendermint}}$
  under bounded delay attack
  ($|V|=4$, $f=1$, $\Delta=\Sigma=1$, $H_{\max}=1$,
  \texttt{length}$=5$).
  RC$\,=\,0$ indicates that no Agreement counterexample was found
  within the bounded computation.
}
\label{tab:safety}
\small
\begin{tabular}{@{}ccrrc@{}}
\toprule
\textbf{Run} &
$R_{\max}$ &
\textbf{length} &
\textbf{Time\,(s)} &
\textbf{Result} \\
\midrule
SR-01 &  1 & 5 & 61.6 & PASS \\
SR-02 &  2 & 5 & 60.6 & PASS \\
SR-03 &  3 & 5 & 57.7 & PASS \\
SR-04 &  5 & 5 & 48.1 & PASS \\
SR-05 &  8 & 5 & 58.3 & PASS \\
SR-06 & 10 & 5 & 44.9 & PASS \\
\bottomrule
\end{tabular}
\end{table}

% ============================================================

% ==================================================================
% 1) TERMINATION
% ==================================================================
\subsection{Termination Verification}

For termination, we only use two key notions:
\(\mathrm{InitiallyHonest}\) and \(\mathrm{NotYetTerminated}\).

\begin{align*}
\mathrm{InitiallyHonest}
&\triangleq
V \setminus V_{\mathrm{corrupted}} \\
\mathrm{NotYetTerminated}
&\triangleq
\neg\Big(\forall p\in\mathrm{InitiallyHonest},\; \mathrm{decision}_p(0)\neq \bot\Big)
\end{align*}

Interpretation:
\(\mathrm{InitiallyHonest}\) is the set of nodes not initially corrupted;
\(\mathrm{NotYetTerminated}\) means that at least one initially honest node
has not decided at height \(0\).
In our runs, a counterexample to \texttt{NotYetTerminated}
(i.e., RC\(=12\)) is treated as a termination witness.

For Table~\ref{tab:termination}, each column has the following meaning:
\textbf{Run} is the experiment identifier,
$R_{\max}$ is the maximum round bound used in that run,
\textbf{length} is the simulation trace-length bound passed to Apalache,
\textbf{Time (s)} is wall-clock verification time in seconds,
and \textbf{Result} indicates whether a termination witness was found
(PASS corresponds to RC$=12$ in this setup).

\begin{table}[ht]
\centering
\caption{%
  Termination verification results for
  $\mathcal{F}^{V,\Delta,\Sigma}_{\mathrm{Tendermint}}$
  under bounded delay attack
  ($|V|=4$, $f=1$, $\Delta=\Sigma=1$, $H_{\max}=1$,
  \texttt{length}$=100$, \texttt{max-run}$=60$).
  RC$\,=\,12$ indicates that Apalache found a termination witness.
}
\label{tab:termination}
\small
\begin{tabular}{@{}ccrrc@{}}
\toprule
\textbf{Run} &
$R_{\max}$ &
\textbf{length} &
\textbf{Time\,(s)} &
\textbf{Result} \\
\midrule
TR-01 &  1 & 100 &    31.3 & PASS \\
TR-02 &  2 & 100 &    25.4 & PASS \\
TR-03 &  3 & 100 &    34.0 & PASS \\
TR-04 &  5 & 100 &   117.1 & PASS \\
TR-05 &  8 & 100 &    81.4 & PASS \\
TR-06 & 10 & 100 &  1097.1 & PASS \\
\bottomrule
\end{tabular}
\end{table}

% ==================================================================
% 2) SAFETY
% ==================================================================
\subsection{Safety Verification}

To verify safety of $\mathcal{F}^{V,\Delta,\Sigma}_{\mathrm{Tendermint}}$
under delay attacks, we use Apalache in bounded checking mode
(\texttt{apalache-mc check}) and evaluate the invariant
$\texttt{Agreement}$, requiring that no two initially honest validators
decide different values at the same height.

The safety property in standard logical form is:
\begin{align*}
\mathrm{Agreement}
&\triangleq
\forall h\in\{0,\dots,H_{\max}\},\; \forall p,q\in\mathrm{InitiallyHonest},
\\
&\qquad\Big(\mathrm{decision}_p(h)\neq\bot \wedge \mathrm{decision}_q(h)\neq\bot\Big)
\Rightarrow
\mathrm{decision}_p(h)=\mathrm{decision}_q(h)
\end{align*}
We fix the validator set $V=\{v_1,v_2,v_3,v_4\}$ with fault threshold $f=1$,
network delay bound $\Delta=1$, protocol attack tolerance $\Sigma=1$,
and block height bound $H_{\max}=1$.

The checker outcome is interpreted as follows:
RC$\,=\,0$ (NoError) means no Agreement counterexample exists within the
bounded horizon (here \texttt{length}$=5$), while RC$\,=\,12$ would indicate
that a violating trace was found.
Table~\ref{tab:safety} shows RC$=0$ for every tested $R_{\max}$,
so no safety violation is observed in this configuration.

Similar to other bounded symbolic checks, runtime is not strictly monotonic
in $R_{\max}$: solver effort depends on the structure of enabled/disabled
transitions and pruning opportunities, not only on the numeric round bound.

For Table~\ref{tab:safety}, the column semantics are:
\textbf{Run} denotes the experiment identifier,
$R_{\max}$ is the maximum round bound for that check,
\textbf{length} is the bounded model-checking horizon,
\textbf{Time (s)} is wall-clock runtime in seconds,
and \textbf{Result} reports the Agreement-check outcome
(PASS corresponds to RC$=0$, i.e., no counterexample found within bounds).

\begin{table}[ht]
\centering
\caption{%
  Safety verification results for
  $\mathcal{F}^{V,\Delta,\Sigma}_{\mathrm{Tendermint}}$
  under bounded delay attack
  ($|V|=4$, $f=1$, $\Delta=\Sigma=1$, $H_{\max}=1$,
  \texttt{length}$=5$).
  RC$\,=\,0$ indicates that no Agreement counterexample was found
  within the bounded computation.
}
\label{tab:safety}
\small
\begin{tabular}{@{}ccrrc@{}}
\toprule
\textbf{Run} &
$R_{\max}$ &
\textbf{length} &
\textbf{Time\,(s)} &
\textbf{Result} \\
\midrule
SR-01 &  1 & 5 & 61.6 & PASS \\
SR-02 &  2 & 5 & 60.6 & PASS \\
SR-03 &  3 & 5 & 57.7 & PASS \\
SR-04 &  5 & 5 & 48.1 & PASS \\
SR-05 &  8 & 5 & 58.3 & PASS \\
SR-06 & 10 & 5 & 44.9 & PASS \\
\bottomrule
\end{tabular}
\end{table}
